\section{Laufzeitanalyse}
\label{sec:runtime_analysis}

Im Folgenden werden die diskrete Convolution in der Zeitdomäne und die Variante der Convolution,
basierend auf der Frequenzdomäne gegenübergestellt.

Dazu betrachten wir zunächst folgenden Pseudocode, der die diskrete Convolution in der
Zeitdomäne algorithmisch aufzeigt:

% Farben für Code definieren
\definecolor{codegreen}{rgb}{0,0.6,0}
\definecolor{codegray}{rgb}{0.5,0.5,0.5}
\definecolor{codepurple}{rgb}{0.58,0,0.82}
\definecolor{backcolour}{rgb}{0.95,0.95,0.92}

% Eigene Sprache für Pseudocode
\lstdefinelanguage{Pseudo}{
    morekeywords={
        algorithm, input, output, function, procedure,
        if, then, else, end, while, do, for, foreach,
        return, break, continue, repeat, until
    },
    sensitive=false,
    morecomment=[l]{//},
    morestring=[b]"
}

% Stil für Pseudocode
\lstdefinestyle{pseudocode}{
    language=Pseudo,
    backgroundcolor=\color{backcolour},
    commentstyle=\color{codegreen},
    keywordstyle=\color{blue},
    numberstyle=\tiny\color{codegray},
    stringstyle=\color{codepurple},
    basicstyle=\ttfamily\small,
    breakatwhitespace=false,
    breaklines=true,
    captionpos=b,
    keepspaces=true,
    numbers=left,
    numbersep=5pt,
    showspaces=false,
    showstringspaces=false,
    showtabs=false,
    tabsize=2,
    mathescape=true
}

\begin{lstlisting}[
    style=pseudocode,
    caption={Grundalgorithmus der Convolution in der Zeitdomäne},
    label={lst:directConvolution}
]
Algorithm directConvolution(inSignal, N, irSignal, M)
    outLen <- N + M - 1     // O(1)
    for i <- 0 to outLen - 1    // O(N + M + 1)
        for k <- 0 to M - 1     // O(M)
            if i-k < 0 or n <= i-k  // O(1)
                continue
            
            outSignal[i] <- outSignal[i]
            + irSignal[k] * inSignal[i-k] // Worst-Case: O(1)
        done
    done
    return outSignal
\end{lstlisting}

Als Parameter erhält der Algorithmus der direkten Convolution zwei Listen.
Die erste Liste enthält die Folge der diskreten Abtastungssignale des Eingangssignals.
Dieses Signal hat den Betrag/ die Länge \textit{N}.
Die zweite Liste enthält die diskreten Abtastungssignale der Impulsantwort und hat die Länge \textit{M}.
Insgesamt hat die resultierende Folge eine Länge von $N + M - 1$.

Über diese Länge iteriert dementsprechend die äußere Schleife.
Die innere Schleife iteriert über die Länge $M$ der Impulsantwort.
Die if-Abfrage der inneren Schleife prüft, ob im Laufe der Iterationen
der Index des Eingangssignals 1 unterschreitet.
Durch diese Abzweigung an Sezenarien ergibt sich eine Worst-Case Laufzeit von
\begin{equation}
\mathcal{O}(1 + (N + M - 1) \cdot (M + 2)) = \mathcal{O}(M^2 + NM - M + 2N - 1)
\end{equation}

Laut Polynomregel fallen Terme niederer Ordnung weg. Daraus folgt:
\begin{equation}
    \mathcal{O}(M^2 + MN) \implies \mathcal{O}(N \cdot M)
     , denn: N > M \implies N \cdot M > M^2
\end{equation}

Die Laufzeitklasse der direkten, diskreten Convolution beträgt also in O-Notation
\begin{equation}
\mathcal{O}(N \cdot M)
\end{equation}

Nun zur Convolution per Multiplikation beider Signale innerhalb der Frequenzdomäne.
Auch hier der Algorithmus in Pseudocode:

\begin{lstlisting}[
    style=pseudocode,
    caption={Grundalgorithmus der Convolution in der Zeitdomäne},
    label={lst:directConvolution}
]
Algorithm fftConvolution(inSignal, N, irSignal, M)
    outLen <- N + M - 1 //  O(1)
    fftLen <- zeroPadToNextPowOf2(outLen)    //  O(1)

    // zeroPad() fuellt die Listen bis fftLen mit nullen auf
    zeroPad(inSignal, fftLen) // O(n)
    zeroPad(irSignal, fftlen) // O(n)

    inSpectrum <- FFT(inSignal, N, fftLen)  // O(n log(n))
    irSpectrum <- FFT(irSignal, M, fftLen)  // O(n log(n))

    for f <- 0 to fftLen - 1    // Insg. O(n)
        outSpectrum[f] <- inSpectrum[f] * irSpectrum[f]
    done

    outTimeDomain <- IFFT(outSpectrum, fftLen)  // O(n log(n))

    // real() isoliert den Realteil des komplexen IFFT-Outputs
    for i <- 0 to outLen - 1    // Insg. O(n)
        outSignal[i] <- real(outTimeDomain[i])
    done

    return outSignal
\end{lstlisting}

Dieser Algorithmus erhält als Parameter erneut die beiden Folgen der diskreten Signale
des Eingangssignals und der Impulsantwort.
Zuerst wird die Länge \textit{fftLen} des Ausgangssignals ermittelt und die beiden Listen mit bis zu
dieser Länge mit Nullen aufgefüllt.
Die Funktion FFT überführt die beiden Listen jeweils in ihre Frequenzdomäne.
Anschließend folgt über die Länge von \textit{fftLen} eine Multipklikation jedes Frequenzwertes.

Die Laufzeitklasse der FFT ist $\mathcal{O}(N \cdot log(N))$ \cite[Kap. 12] {smith1997}. Durch Addition aller anderen
Operationen ergibt sich eine Laufzeit von:
\begin{equation}
    \mathcal{O}(3(N \cdot log(N)) + 3N + 2) \in \mathcal{O}(N \cdot log(N))
\end{equation}